% Generated draft: requires mathematical and source review.
\documentclass[11pt]{article}
\usepackage[margin=25mm]{geometry}
\usepackage{fontspec}
\setmainfont{Latin Modern Roman}
\usepackage{amsmath,amssymb,mathtools,unicode-math}
\setmathfont{Latin Modern Math}
\usepackage{enumitem,needspace,xurl,hyperref}
\setlength{\parindent}{0pt}
\setlength{\parskip}{5pt}
\setlength{\emergencystretch}{3em}
\newcommand{\N}{\mathbb{N}}
\newcommand{\Z}{\mathbb{Z}}
\newcommand{\Q}{\mathbb{Q}}
\newcommand{\R}{\mathbb{R}}
\newcommand{\C}{\mathbb{C}}
\newcommand{\F}{\mathbb{F}}
\newcommand{\A}{\mathbb{A}}
\newcommand{\PP}{\mathbb{P}}
\newcommand{\E}{\mathbb{E}}
\newcommand{\eps}{\varepsilon}
\newcommand{\dd}{\,\mathrm{d}}
\begin{document}
\section*{UnsolvedMath ID: 30005310}
\textbf{Source problem code:} catalogue entry
\subsection{Definitions and mathematical statement}
Let \(G\) be a colored Gaussian graphical model, and let \(n\) be the rank index appearing in the supplied formulation. Assume
\[
I(G,n)\ne(0),\qquad I(G,n+1)=(0),
\]
where the record calls \(I(G,n)\) the rank-\(n\) elimination ideal. The construction of this ideal, its coefficient field, and its ambient polynomial ring are not supplied.

Let \(p\) denote the number of Gaussian variables, let \(\operatorname{Sym}_p(\R)\) denote the real symmetric \(p\times p\) matrices, and write
\[
\operatorname{Sym}_p^{++}(\R)
=\{\Sigma\in\operatorname{Sym}_p(\R):x^{\mathsf T}\Sigma x>0
\text{ for every }x\in\R^p\setminus\{0\}\}.
\]
Write \(T_G(\Sigma)\) for the sufficient-statistic vector associated with \(\Sigma\). Its coordinates and normalization are not specified in the record. If \(t=(t_1,\ldots,t_m)\) are the intended sufficient-statistic coordinates, the real zero set relevant to the covariance condition is
\[
V_{\R}(I(G,n))
=\{t\in\R^m:f(t)=0\text{ for every }f\in I(G,n)\}.
\]
This interpretation requires the omitted identification of the ideal's variables with the sufficient-statistic coordinates; the record itself writes \(V(I(G,n))\) without specifying the base field.

Denote by \(\mathcal A(G,n)\) the source's condition that there is a ``nonvanishing signed sum of squares in \(I(G,n)\).'' This is a placeholder for an incompletely specified condition: the permissible signs, the variables in which squares are taken, any substitution, and the domain on which nonvanishing is required are not supplied. Define the covariance condition, subject to the coordinate qualification above, by
\[
\mathcal B(G,n)\quad\Longleftrightarrow\quad
\exists\Sigma\in\operatorname{Sym}_p^{++}(\R)
\text{ such that }T_G(\Sigma)\in V_{\R}(I(G,n)).
\]
The supplied formulation refers to a positive-definite covariance matrix without stating an additional model-membership restriction.

\textbf{Conjecture 4, as reconstructed in the supplied registry.}
For every such \(G\) and rank index \(n\), exactly one of the following scenarios occurs:
\begin{enumerate}
\item \(\mathcal A(G,n)\) holds and
\[
\operatorname{WMLT}(G)=\operatorname{MLT}(G)=n+1.
\]
\item \(\mathcal B(G,n)\) holds and
\[
\operatorname{WMLT}(G)=n<\operatorname{MLT}(G)=n+1.
\]
\end{enumerate}
Here \(\operatorname{WMLT}\) and \(\operatorname{MLT}\) are the weak maximum-likelihood threshold and maximum-likelihood threshold named in the record. Their sampling conventions, probability quantifiers, and boundary conventions are absent. Consequently this is an explicitly qualified statement of the intended question, not a complete self-contained formulation of Proposition 3's two cases.

The original excerpt states only that the two cases in Proposition 3 are the only possible scenarios, and then includes surrounding example material:
\[
\begin{aligned}
&4(-l_{14}l_{23}+l_{13}l_{24})^2
+4(-l_{14}l_{33}+l_{13}l_{34})^2
+4(l_{13}l_{44})^2\\
&\quad+4(l_{22}l_{33})^2
+4(-l_{24}l_{33}+l_{23}l_{34})^2
+4(l_{23}l_{44})^2
+4(l_{33}l_{44})^2.
\end{aligned}
\]
The excerpt describes this expression as a nonvanishing sum of squares because \(l_{11},l_{22},l_{33},l_{44}>0\), and reports \(\operatorname{WMLT}(G_9)=\operatorname{MLT}(G_9)=2\). Neither \(G_9\) nor the variables \(l_{ij}\), their substitution, or their relation to the elimination ideal is defined in the supplied excerpt. This example does not supply the missing statement of Proposition 3.
\subsection{Short English statement}
The supplied conjecture asks whether two mechanisms exhaust the threshold behavior of colored Gaussian graphical models under the stated elimination-ideal assumptions: a nonvanishing signed sum-of-squares condition giving both thresholds equal to \(n+1\), or a positive-definite covariance condition giving weak threshold \(n\) and maximum-likelihood threshold \(n+1\). Definitions needed to make these alternatives fully precise are missing, particularly Proposition 3 and the threshold conventions.
\subsection{Sources}
\begin{itemize}
\item Supplied local UnsolvedMath registry record, permanent ID \(30005310\), title: \emph{Exhaustiveness of Threshold Scenarios for Colored Gaussian Graphical Models}; exact source problem code: \texttt{catalogue entry}. The record attributes the question to Conjecture 4 in an Oberwolfach workshop report; no individual proposer is supplied.
\item Supplied source citation: \emph{Algebraic Structures in Statistical Methodology} (2023), Oberwolfach Reports, pp.\ 3121--3170, DOI \texttt{10.4171/owr/2022/55}. Original report link: \url{https://doi.org/10.4171/owr/2022/55}. The registry gives proposed year 2022; both this year and the supplied 2023 citation are retained without resolving their bibliographic relationship.
\item Supplied literature source labeled ``Bernstein et al. primary article'': \url{https://arxiv.org/abs/2108.02185}. The record does not provide its title or a precise theorem connecting it to the conjecture.
\item Literature assessment dated exactly \texttt{2026-08-22}: the supplied curation labels the problem open and says that the exact two-scenario claim remains a conjecture, while rigidity-based maximum-likelihood-threshold results provide substantial structure without establishing the exhaustive classification. This is supplied dated curation, not an independently verified literature assessment. The background says its sources checked were the original OWR report and the Bernstein et al. article linked above, and requests independent verification of time-sensitive status and cited claims.
\item The registry labels the statement \texttt{corrected\_verified} and says the official report was checked and its notation repaired. These are supplied provenance claims; no independent source verification is asserted in this draft.
\end{itemize}
\end{document}
